\documentclass[12pt]{article}
\usepackage[utf8]{inputenc}
\usepackage[T1]{fontenc}
\usepackage{hyperref}
\title{Tututni Language Resources for K-12 Curriculum Development}
\author{Compiled by Claude Sonnet 5 (claude-sonnet-5) for Andr\'as Kornai}
\date{}
\begin{document}
\maketitle

\section{Introduction}
Tututni (ISO 639-3 tuu) is an Athabaskan language of the southern Oregon coast (Rogue River area). It may have no remaining native speakers. Search found very little: a Smithsonian note about a historical manuscript and general reference pages. No modern teaching materials or school programme were found. Related Coquille and Siletz tribal programmes exist, but they are other varieties or languages and should not be confused with Tututni.

\section{Primary Educational Resources}
None found.

\section{Linguistic and Reference Materials}
\begin{itemize}
\item Smithsonian Recovering Voices: ``Can a document help revitalize the Tututni language?'' This describes a roughly 200-page vocabulary manuscript compiled by James Owen Dorsey in the 1880s, held at the National Anthropological Archives. \url{https://naturalhistory.si.edu/research/anthropology/programs/recovering-voices/outreach/news-research-updates/can-document-help-revitalize-tututni-language}
\item Tututni overview in the Athabaskan documentation materials at the University of Missouri. I saw only the title and URL in search results and did not check the content. \url{https://dice.missouri.edu/assets/docs/nadene/Tutuni.pdf}
\item Wikipedia, Coquille people, which mentions Coquille revitalization work on Upper Coquille Athabaskan. \url{https://en.wikipedia.org/wiki/Coquille_people}
\end{itemize}

\section{Religious and Community Resources}
The Coquille Indian Tribe is reported to run revitalization work on Upper Coquille Athabaskan, a related Lower Rogue River variety (per Wikipedia). The Confederated Tribes of Siletz run a programme for Siletz Dee-ni, which is a different language. No Tututni-specific community organisation was found.

\section{Digital and Audio Resources}
None found.

\section{Curriculum Development Context}
Tututni is not known to be taught in schools. Any curriculum would be a heritage or revitalization effort built on archival sources such as the Dorsey manuscript, and it should involve descendant communities (for example Coquille and Siletz) from the start. Materials from related Athabaskan languages might be adapted, but this needs community consent and linguistic review.
\end{document}